\chapter{$\mathrm{GL}_n$ 的 Langlands 猜想}\label{ch:10}

\begin{center}
J.\,W. Cogdell
\end{center}

现代数论的主要目标之一，是理解局部域或整体域 $k$（例如 $\mathbb{Q}$）的 Galois 群
$\mathcal{G}_k=\operatorname{Gal}(\bar k/k)$。理解群 $\mathcal{G}_k$ 的一条途径，
是理解它的有限维表示论。当 $k$ 为数域时，Artin 对每个有限维表示
$\rho\colon\mathcal{G}_k\to\GL_n(\mathbb{C})$ 附上一个复解析不变量，即其 $L$ 函数
$L(s,\rho)$。于是理解 $\rho$ 的一条途径就是通过这个不变量。当 $\rho$ 是一维表示时，
这一想法是 Abel 类域论的解析进路以及理解 $\mathcal{G}_k^{\mathrm{ab}}$ 的基础。
为了更完整地理解 $\mathcal{G}_k$，我们自然希望更完整地理解高维表示的 $L(s,\rho)$。
\begin{annote}
一维情形是全部类域论的解析核心。局部类域论给出 $k^\times$ 的乘法特征与
$\mathcal{G}_k^{\mathrm{ab}}$ 的连续特征之间的典范双射；整体情形则给出
$k^\times\backslash\mathbb{A}^\times$ 的 Hecke 特征与 $\mathcal{G}_k^{\mathrm{ab}}$
的特征之间的双射。正文所说的「解析进路」指 Tate 的博士论文：把 Hecke 特征写成 Euler
乘积，再用局部 $\zeta$ 积分证明函数方程。另需注意，本章自始至终采用几何 Frobenius 的
归一化，故 $\omega_\pi$ 与 $\det\rho_\pi$ 的对应不差 $q$ 的幂；若换用算术 Frobenius，
所有 $\varepsilon$ 因子都要整体扭转一个 $\omega^s$ 型的特征。
\par\vspace{0.15em}{\raggedleft 【来源：Tate, \textit{Fourier analysis in number fields and Hecke's zeta-functions}（1950 年学位论文，收入 Cassels--Fr\"ohlich 编 \textit{Algebraic Number Theory}, Academic Press 1967, 305--347）；Tate, \textit{Number theoretic background}, Proc.\ Symp.\ Pure Math.\ 33 (1979), Part 2, 3--26（A）】\par}
\end{annote}

另一类对象也具有类似的解析不变量，即 $\GL_n(\mathbb{A})$ 的自守表示 $\pi$，
其中 $\mathbb{A}$ 是 $k$ 的阿代尔环。自守表示所附的 $L$ 函数 $L(s,\pi)$
的解析性质已有很好的理解 [13]。

Langlands 猜想预言：$\mathcal{G}_k$ 的 $n$ 维表示与 $\GL_n(\mathbb{A})$ 的自守表示之间
存在一种对应，并且保持这些解析不变量。相应地，局部域 $k$ 的 $\mathcal{G}_k$ 的 $n$ 维表示
与 $\GL_n(k)$ 的容许表示之间也有一种对应，即局部 Langlands 猜想。看待这类对应有两种方式。
若把信息的传递看作从自守侧流向 Galois 侧（正如上文所做），这就是局部或整体非 Abel 类域论。
若把信息的传递看作从 Galois 侧流向自守侧，这就是容许表示或自守表示的算术参数化。

过去十年间，人们对这些 Langlands 猜想的理解取得了可观的进展。进程始于九十年代初：
Laumon、Rapoport 与 Stuhler 证明了特征 $p$ 局部域 $k$ 的局部 Langlands 猜想 [44]。
九十年代后期，Harris 与 Taylor 证明了特征零非阿基米德域的局部 Langlands 猜想 [27]，
Henniart 随即给出一个简化证明 [30]。大约同时，沿着 Drinfeld 证明整体函数域上
$\GL_2$ 的 Langlands 猜想 [23] 时开创的纲领，L.~Lafforgue 在函数域情形建立了
$\GL_n$ 的整体 Langlands 猜想 [35]。

本综述意在概述这些结果，并着重于它们的共同特征。关于各项独立工作，已有若干出色的综述：
Carayol 关于局部猜想的综述 [9, 10]，以及 Laumon 关于整体猜想的综述 [42, 43]；
更深入的论述可参阅这些文献。本文前两节讨论 Galois 表示、自守表示及其 $L$ 函数。
接着讨论局部 Langlands 猜想的两种形式：表示论形式，由 Langlands 于 1973 年前后
在阿基米德情形证明 [38]；以及 $L$ 函数形式，非阿基米德情形下由 Laumon、Rapoport、
Stuhler、Harris、Taylor 与 Henniart 所建立的形式。最后讨论 Drinfeld 与 Lafforgue
在特征 $p$ 下建立的整体 Langlands 猜想。

尽管数域上的整体 Langlands 猜想鲜有一般性进展，近年来却确立起若干引人注目的特殊情形。
其中最著名的是 Wiles 对 $\mathbb{Q}$ 上椭圆曲线所对应的 $\mathcal{G}_{\mathbb{Q}}$
的某些二维 $\ell$ 进表示的模性的证明，这是他在证明 Fermat 大定理途中得到的 [59]，
此外还有一批相关结果。遗憾的是，本文不讨论这些结果。

\section{Galois 表示及其 \texorpdfstring{$L$}{L} 函数}

设 $G$ 为拓扑群，$F$ 为拓扑域。以 $\operatorname{Rep}_n(G;F)$ 记连续表示
$\rho\colon G\to\GL_n(F)$ 的等价类之集，以 $\operatorname{Rep}_n^0(G;F)$
记其中不可约表示的子集。我们主要关心复表示，故把 $\operatorname{Rep}_n(G;\mathbb{C})$
简记为 $\operatorname{Rep}_n(G)$，$\operatorname{Rep}_n^0$ 同理。有时我们关心
$F=\Ql$ 的情形，此时在记号中保留系数域。

设 $k$ 为局部域或整体域，以 $\bar k$ 记 $k$ 的可分代数闭包。令
$\mathcal{G}_k=\operatorname{Gal}(\bar k/k)$ 为（绝对）Galois 群，$W_k$ 为（绝对）
Weil 群 [51]。

设 $k$ 为非阿基米德局部域，特征为 $p$，剩余域 $\kappa$ 的阶为 $q$。
令 $I\subset\mathcal{G}_k$ 为惯性子群。取 $\Phi$ 为 $\mathcal{G}_k$ 的一个几何
Frobenius 元，则 $W_k$ 可取为 $\mathcal{G}_k$ 中由 $\Phi$ 与 $I$ 代数生成的子群，
其拓扑使得 $I$ 带有从 $\mathcal{G}_k$ 诱导的拓扑、$I$ 为开子群、且左乘 $\Phi$ 为同胚。
它还可被赋予 $\mathbb{Q}$ 上的概形结构 [51]。于是有连续同态
$W_k\to\mathcal{G}_k$，其像在 $\mathcal{G}_k$ 中稠密。因此有自然嵌入
$\operatorname{Rep}_n(\mathcal{G}_k)\hookrightarrow\operatorname{Rep}_n(W_k)$。
其像，即那些可分解为 $\mathcal{G}_k$ 的连续表示的表示，称为 $W_k$ 的 Galois 型表示。
我们还有自然特征 $\omega^s\in\operatorname{Rep}_1(W_k)$，由 $\omega^s(I)=1$ 与
$\omega^s(\Phi)=q^{-s}$ 定义。它也可写作 $\omega^s(w)=\|w\|^{-s}$，并给出同态
$v\colon W_k\to\mathbb{Z}$，其中 $\|w\|=q^{-v(w)}$。于是 $W_k$ 的每个不可约表示
$\rho$ 都有形状 $\rho=\rho^0\otimes\omega^s$，其中 $\rho^0$ 为 Galois 型 [49, 19]。

在算术代数几何中最自然出现的表示，例如与代数簇的 $\ell$ 进上同调相联系的表示，
并非复表示，而是 $\operatorname{Rep}_n(\mathcal{G}_k;\Ql)$（$\ell\neq p$）
或 $\operatorname{Rep}_n(W_k;\Ql)$ 中的表示。由于两个域上拓扑不同，$\ell$ 进表示的
表示论比复表示更为丰富。Deligne 注意到这一点，引入了如今称为局部域 Weil--Deligne 群
$W'_k$ 的对象；其表示论本质上是代数的，因而实质上不区分 $\mathbb{C}$ 与 $\Ql$，
其表示范畴与 $\mathcal{G}_k$ 或 $W_k$ 的连续 $\ell$ 进表示范畴相同 [19]。
仿照 Tate [51]，我们如下定义 $W'_k$：它是 $\mathbb{Q}$ 上的群概形，即 Weil 群 $W_k$
与加法群 $\mathbb{G}_a$ 的半直积，$W'_k=W_k\ltimes\mathbb{G}_a$，其中 $W_k$ 通过
$wxw^{-1}=\|w\|x$ 作用在 $\mathbb{G}_a$ 上。若 $F$ 为任意特征零的域，例如 $\Ql$
或 $\mathbb{C}$，则 $W'_k$ 的 $F$ 点就是 $W_k\times F$，其乘法为
\[
  (w_1,x_1)(w_2,x_2)=(w_1w_2,\;x_1+\|w_1\|x_2).
\]
但真正重要的是 $W'_k$ 的表示论。$W'_k$ 在 $F$ 上的 $n$ 维表示是一对 $\rho'=(\rho,N)$，
包括：(i)~$n$ 维 $F$ 向量空间 $V$ 与群同态 $\rho\colon W_k\to\GL(V)$，其核包含 $I$ 的
开子群，也就是说，$\rho$ 关于 $\GL(V)$ 上的离散拓扑连续；(ii)~$V$ 的幂零自同态 $N$，
满足 $\rho(w)N\rho(w)^{-1}=\|w\|N$ [19, 49, 51]。
\begin{annote}
半直积的乘法里，交叉项必须写成 $x_1+\|w_1\|x_2$。验算结合律：
\[
  \bigl((w_1,x_1)(w_2,x_2)\bigr)(w_3,x_3)
  =(w_1w_2w_3,\;x_1+\|w_1\|x_2+\|w_1w_2\|x_3),
\]
与先乘后两项所得一致，关键在于 $\|w_1w_2\|=\|w_1\|\,\|w_2\|$。若误写成
$x_1+\|w_2\|x_2$，结合律立刻失效。另需注意，$W'_k$ 的「表示」按定义是 $(\rho,N)$
而非群同态：$N$ 记录的是 $\mathbb{G}_a$ 在李代数层面的作用，不能丢弃。
\par\vspace{0.15em}{\raggedleft 【来源：Deligne, \textit{Les constantes des \'equations fonctionnelles des fonctions $L$}, Lecture Notes in Math.\ 349 (1973), 501--597；Tate 1979 同上（A）】\par}
\end{annote}

设 $\rho'=(\rho,N)$ 为 $W'_k$ 的表示。则存在 $V$ 的唯一幂幺自同构 $u$，
它与 $N$ 及 $\rho(W_k)$ 都可交换，且使 $e^{aN}\rho(w)u^{-v(w)}$ 对一切 $a\in F$、
$w\in W_k-I$ 半单 [19, 51]。$\rho'$ 的 $\Phi$ 半单化即
$\rho'_{\mathrm{ss}}=(\rho u^{-v},N)$。若 $\rho'=\rho'_{\mathrm{ss}}$，则称 $\rho'$
为 $\Phi$ 半单的（或 Frobenius 半单的）；此时 $u$ 为恒等，所有 Frobenius 元都半单地作用。
这等价于表示 $\rho$ 在通常意义下半单。

以 $\operatorname{Rep}_n(W'_k;F)$ 记 Weil--Deligne 群 $W'_k$ 的 $n$ 维 $\Phi$ 半单
$F$ 表示的等价类之集。当 $F=\mathbb{C}$ 时，把
$\operatorname{Rep}_n(W'_k;\mathbb{C})$ 简记为 $\operatorname{Rep}_n(W'_k)$。

Weil--Deligne 群的重要性在于：它让我们以代数的方式捕捉 $\mathcal{G}_k$ 或 $W_k$
的连续 $\ell$ 进表示 [19, 49, 51]。对每个半单 $\ell$ 进表示
$\rho_\ell\in\operatorname{Rep}_n(W_k;\Ql)$，惯性群 $I$ 中存在开子群使 $\rho_\ell$
在其上平凡，于是 $\rho_\ell$ 给出 $W'_k$ 的一个（通常意义下的）$\Phi$ 半单 $\Ql$ 表示
$\rho'$。注意，由 $W'_k$ 表示定义中的条件 (ii)，$F$ 上的拓扑不起作用；故若固定同构
$\iota\colon\Ql\to\mathbb{C}$，便可等同
$\operatorname{Rep}_n(W'_k;\Ql)\cong\operatorname{Rep}_n(W'_k;\mathbb{C})
=\operatorname{Rep}_n(W'_k)$。此外，$W'_k$ 的不可约表示必有 $N=0$，因为 $N$ 的核
将是子表示；于是 $\operatorname{Rep}_n^0(W'_k)=\operatorname{Rep}_n^0(W_k)$。

设 $\rho'=(\rho,N)\in\operatorname{Rep}_n(W'_k;F)$ 是 $W'_k$ 在向量空间 $V$ 上的表示。
令 $V_N^I=(\ker N)^{\rho(I)}$ 为惯性子群 $I$ 在 $\ker N$ 上的不变子空间。于是可如下
定义局部 $L$ 因子：
\[
  Z(t,V)=\det\bigl(1-t\rho(\Phi)|_{V_N^I}\bigr)^{-1}\in F(t),
\]
即 $\Phi$ 在 $V_N^I$ 上作用的特征多项式的倒数。若有嵌入 $F\hookrightarrow\mathbb{C}$
（例如 $F=\mathbb{C}$，或利用同构 $\iota\colon\Ql\to\mathbb{C}$），就把 $F$ 看作
$\mathbb{C}$ 的子域，并令
\[
  L(s,\rho')=Z(q^{-s},\rho').
\]
\begin{annote}
局部 $L$ 因子只在 $\ker N$ 的惯性不变向量上取 $\Phi$ 的特征多项式。$N\neq0$ 时
$V^I$ 会偏大：以 $\GL_2$ 的 Steinberg 表示为例，对应的
$\rho'=\omega\otimes\mathrm{Sp}(2)$ 有 $N\neq0$，若改用 $V^I$ 就会把 $L$ 因子多算
一个因子。对未分歧表示则 $N=0$ 且 $I$ 平凡地作用，$V_N^I=V$，于是回到熟悉的
$\det(1-\rho(\Phi)q^{-s})^{-1}$。
\par\vspace{0.15em}{\raggedleft 【来源：Tate 1979, \S4；Rohrlich, \textit{Elliptic curves and the Weil--Deligne group}, CRM Proc.\ Lecture Notes 4 (1994), 125--157（A）】\par}
\end{annote}

局部常数 $\varepsilon(s,\rho',\psi)$（其中 $\psi$ 为 $k$ 的加法特征）的定义更为精细，
其确切定义可参看 Deligne [19]、Rohrlich [49] 或 Tate [51]。当然，当 $N=0$ 时，
这些就是通常的局部 Artin--Weil $L$ 函数与 $\varepsilon$ 因子。

若局部域 $k$ 是阿基米德的，即 $k=\mathbb{R}$ 或 $\mathbb{C}$，则我们只关心
$\mathcal{G}_k$ 或 $W_k$ 的复表示。当 $k=\mathbb{C}$ 时，Weil 群就是 $\mathbb{C}^\times$；
当 $k=\mathbb{R}$ 时，$W_{\mathbb{R}}\cong\mathbb{C}^\times\sqcup j\mathbb{C}^\times$，
其中 $j^2=-1$，且对 $c\in\mathbb{C}^\times$ 有 $jcj^{-1}=\bar c$。两种情形都有正合列
\[
  1\longrightarrow\mathbb{C}^\times\longrightarrow W_k
  \longrightarrow\mathcal{G}_k\longrightarrow 1.
\]
不存在阿基米德 Weil--Deligne 群，为统一起见，在这些情形令 $W'_k=W_k$。
$L$ 因子与 $\varepsilon$ 因子则用经典 $\Gamma$ 函数与一个局部函数方程来定义 [51]。
\begin{annote}
$W_{\mathbb{R}}$ 作为集合是 $\mathbb{C}^\times$ 与 $j\mathbb{C}^\times$ 的无交并，
但它不是 $\mathbb{C}^\times$ 与 $\mathbb{Z}/2\mathbb{Z}$ 的直积：$jzj^{-1}=\bar z$
给出非平凡的作用，而 $j^2=-1$ 说明正合列
$1\to\mathbb{C}^\times\to W_{\mathbb{R}}\to\mathcal{G}_{\mathbb{R}}\to1$
不分裂。$W_{\mathbb{C}}=\mathbb{C}^\times$ 交换，$W_{\mathbb{R}}$ 非交换；这一点在
比较实群表示与复群表示时必须记住。阿基米德情形没有 $W'_k$，故 $L$ 与 $\varepsilon$
因子只能用经典 $\Gamma$ 函数与局部函数方程定义，不能像非阿基米德情形那样借 $N$ 的核
来截断。
\par\vspace{0.15em}{\raggedleft 【来源：Tate 1979, \S1--2（A）】\par}
\end{annote}

当 $k$ 为整体域时，我们至少关心整体 Galois 群 $\mathcal{G}_k$、整体 Weil 群 $W_k$，
或可能的猜想中的 Langlands 群 $\mathcal{L}_k$ [46] 的表示。

当 $F=\Ql$ 时，以 $\operatorname{Rep}_n(G;\Ql)$ 记如下意义下的整体 $\ell$ 进表示之集：
它们应当连续、代数（即取值于 $\GL(E_\lambda)$，其中 $E_\lambda/\mathbb{Q}_\ell$
为有限扩张），并且几乎处处无分歧（即存在 $k$ 的有限位点集 $S(\rho)$，使得对一切
$v\notin S(\rho)$，表示 $\rho$ 在 $v$ 处无分歧）。

对 $\mathcal{G}_k$ 或 $W_k$ 的任意整体表示 $\rho$，把 $\rho$ 与自然映射
$\mathcal{G}_{k_v}\to\mathcal{G}_k$ 或 $W_{k_v}\to W_k$ 复合，便得到 $k$ 的每个
完备化 $v$ 处的局部表示 $\rho_v$。猜想中的 Langlands 群 $\mathcal{L}_k$ 应当与局部
Weil--Deligne 群具有类似的局部--整体相容性。

对 Galois 群或 Weil 群的任一 $n$ 维复表示或 $\ell$ 进表示，我们都附上一个整体复解析
不变量，即由 Euler 乘积定义的整体 $L$ 函数 $L(\rho,s)$：
\[
  L(s,\rho)=\prod_v L(s,\rho_v),\qquad
  \varepsilon(s,\rho)=\prod_v\varepsilon(s,\rho_v,\psi_v),
\]
其中 $\psi=\prod_v\psi_v$ 为 $k$ 的加法特征。

这些整体解析不变量被猜想具有良好的性质，即：
\begin{enumerate}[label=(\arabic*),leftmargin=*]
  \item $L(s,\rho)$ 应有亚纯延拓，至多有限多个极点；若 $\rho$ 不可约且非平凡，则整；
  \item 这些延拓在竖直带内应有界；
  \item 它们满足函数方程
    $L(s,\rho)=\varepsilon(s,\rho)L(1-s,\tilde\rho)$。
\end{enumerate}
\begin{annote}
$\tilde\rho$ 与 $\tilde\pi$ 是对偶表示（contragredient），不是复共轭表示。对 $\ell$
进表示，须先在 $\Ql$ 上取对偶，再经 $\iota$ 与复表示比较；两者只有在酉情形才与复共轭
一致。函数方程右侧必须是 $\tilde\rho$：写成 $L(1-s,\rho)$ 一般不成立，这是最常见的
笔误之一。
\par\vspace{0.15em}{\raggedleft 【来源：Tate 1979, \S3--4；Cogdell, 本书第 9 章（A）】\par}
\end{annote}

当 $G=\mathcal{G}_k$ 或 $W_k$、$F=\mathbb{C}$ 时，这些就是经典的 Artin--Weil
$L$ 函数；由 Brauer 定理，已知它们在某个右半平面收敛、可亚纯延拓到 $\mathbb{C}$，
并满足函数方程。

若 $k$ 为整体函数域，常数域的阶为 $q$，而 $\rho$ 是如上 $\mathcal{G}_k$ 的 $\ell$
进表示，则 Grothendieck 已证明该 $L$ 函数其实是 $q^{-s}$ 的有理函数，并满足函数方程；
Deligne 随后证明函数方程中的 $\varepsilon$ 因子具有局部分解，且由上述方式给出 [51]。

\section{自守表示及其 \texorpdfstring{$L$}{L} 函数}

在自守侧，若 $k$ 为局部域，以 $\mathcal{A}_n(k)$ 记 $\GL_n(k)$ 的不可约容许复表示的
等价类之集。当 $k$ 为非阿基米德局部域时，以 $\mathcal{A}_n^0(k)$ 记 $\GL_n(k)$ 的
超尖点表示的等价类之子集。由 Godement--Jacquet 理论 [24]，或由 [13] 中概述的
Jacquet--Piatetski-Shapiro--Shalika 理论，每个 $\pi\in\mathcal{A}_n(k)$ 都附有一个
复解析不变量，即其 $L$ 函数 $L(s,\pi)$，以及一个依赖加法特征选取的局部 $\varepsilon$
因子 $\varepsilon(s,\pi,\psi)$。若另有 $\GL_m(k)$ 的不可约容许表示 $\pi'$，
则有局部 Rankin--Selberg 卷积 $L$ 函数 $L(s,\pi\times\pi')$ 与 $\varepsilon$ 因子
$\varepsilon(s,\pi\times\pi',\psi)$。
\begin{annote}
$\mathcal{A}_n(k)$ 在本章有两重含义：$k$ 为局部域时指 $\GL_n(k)$ 的不可约容许表示，
$k$ 为整体域时指 $\GL_n(\mathbb{A})$ 的不可约自守表示。上标 $0$ 在局部指超尖点，
在整体指尖点。正文正是靠这一记号的转换来区分局部与整体两种对应，读者须按所在段落判断，
不可混用。
\par\vspace{0.15em}{\raggedleft 【来源：Cogdell, 本书第 9 章《$\GL_n$ 的 $L$ 函数解析理论》\S1--2（A）】\par}
\end{annote}

若 $k$ 为整体域，以 $\mathbb{A}$ 记其阿代尔环。以 $\mathcal{A}_n(k)$ 记
$\GL_n(\mathbb{A})$ 的不可约自守表示之集，以 $\mathcal{A}_n^0(k)$ 记尖点自守表示
之子集。若 $\pi=\bigotimes'_v\pi_v$ 是 $\GL_n(\mathbb{A})$ 的自守表示，
$\pi'=\bigotimes'_v\pi'_v$ 是 $\GL_m(\mathbb{A})$ 的自守表示，则其附带的 $L$ 函数与
$\varepsilon$ 因子由 Euler 乘积定义：
\[
  L(s,\pi\times\pi')=\prod_v L(s,\pi_v\times\pi'_v),\qquad
  \varepsilon(s,\pi\times\pi')=\prod_v\varepsilon(s,\pi_v\times\pi'_v,\psi_v).
\]

如 [13] 所述，已知这些不变量具有良好性质：若 $\pi$ 与 $\pi'$ 为酉尖点表示，则
\begin{enumerate}[label=(\arabic*),leftmargin=*]
  \item $L(s,\pi\times\pi')$ 可解析延拓到整个 $\mathbb{C}$，且仅当
    $\pi'=\tilde\pi$ 时在 $s=0,1$ 处至多有单极点；
  \item 这些延拓在竖直带内有界；
  \item 满足函数方程
    $L(s,\pi\times\pi')=\varepsilon(s,\pi\times\pi')
    L(1-s,\tilde\pi\times\tilde\pi')$。
\end{enumerate}

考虑出现在 $\ell$ 进上同调中的表示时，最自然的是使用 $\Ql$ 值的自守形式与表示，
记作 $\mathcal{A}_n(k;\Ql)$。例如，我们将需要考虑中心特征为有限阶的 $\Ql$ 值尖点
表示空间，记作 $\mathcal{A}_n^0(k;\Ql)_f$。它们是 $\GL_n(\mathbb{A})$ 的表示，
或等价地，是 $\GL_n(\mathbb{A})$ 上紧支集局部常值 $\Ql$ 值函数所成 Hecke 代数
$\mathcal{H}$ 在 $\GL_n(\mathbb{A})$ 上某些 $\Ql$ 值尖点形式空间中的表示。为便于读者，
我们按 [42] 回顾有限域上函数域情形的定义。我们关心的 $\GL_n(\mathbb{A})$ 上
$\Ql$ 值尖点形式是函数 $\varphi\colon\GL_n(\mathbb{A})\to\Ql$，满足：
\begin{enumerate}[label=(\roman*),leftmargin=*]
  \item 对一切 $\gamma\in\GL_n(k)$ 有 $\varphi(\gamma g)=\varphi(g)$；
  \item 存在紧开子群 $K_\varphi\subset K=\GL_n(\mathcal{O}_k)$，使对一切
    $k\in K_\varphi$ 有 $\varphi(gk)=\varphi(g)$；
  \item 存在 $a\in\mathbb{A}^\times$，$\deg(a)\neq0$，使 $\varphi(ag)=\varphi(g)$；
  \item $\varphi$ 按通常意义是尖点的，即对 $\GL_n(\mathbb{A})$ 的任一极大抛物子群的
    幂幺根 $U$，都有
    \[
      \int_U\varphi(ug)\,du=0.
    \]
\end{enumerate}
\begin{annote}
条件 (iii) 是整体函数域特有的替代品：它要求 $\varphi$ 在某个正次数元素 $a$ 的平移下
不变。由于 $\mathbb{A}^\times/k^\times\cong\mathbb{Z}\times(\text{紧群})$，而
$\langle a\rangle$ 在 $\mathbb{Z}$ 分量中指数有限，中心特征 $\omega_\varphi$ 必在
一个有限指数子群上平凡，故为有限阶。反过来，中心特征有限阶时也可取到这样的 $a$。
这就是正文断言 (iii) 蕴涵中心特征有限阶的理由。
\par\vspace{0.15em}{\raggedleft 【来源：Lafforgue, \textit{Chtoucas de Drinfeld et correspondance de Langlands}, Invent.\ math.\ 147 (2002), 1--241，附录；Laumon, S\'eminaire Bourbaki 873 (2000)（A）】\par}
\end{annote}

注意条件 (iii) 蕴涵 $\varphi$ 的中心特征为有限阶。借助同构
$\iota\colon\Ql\to\mathbb{C}$，这套理论本质上可与复理论等同；相应的 $L$ 函数既可
与通常的复解析 $L$ 函数等同，也可如 [35] 附录那样保留为 $\ell$ 进值有理函数。

\section{局部 Langlands 猜想}

局部 Langlands 猜想最基本的形式，是（Abel）局部类域论的非 Abel 推广。
Langlands 最初表述该猜想时用的是 Weil 群。一个早期的表述（也许是最早的）见于 [36]。
Langlands 从未把自己限制在 $\GL_n$，而总是用一般的约化代数群来表述。Deligne 首先指出，
为了把局部域 $k$ 上 $\GL_2(k)$ 的特殊表示纳入进来，必须过渡到如今所称的
Weil--Deligne 群 [18]。目前最接近 Langlands 原初表述的形式，见于 Borel 在 Corvallis
会议文集中的文章 [4]。

\par\medskip\noindent\textbf{局部 Langlands 猜想 I.}\ 设 $k$ 为局部域。则存在一族自然的双射
\[
  \operatorname{Rep}_n(W'_k)\;\longleftrightarrow\;\mathcal{A}_n(k),
  \qquad \rho=\rho_\pi\;\longleftrightarrow\;\pi=\pi_\rho,
\]
满足一组表示论上的预期性质，其中包括：
\begin{enumerate}[label=(\roman*),leftmargin=*]
  \item 当 $n=1$ 时，它应由局部类域论的同构给出。
  \item $\pi_\rho$ 的中心特征对应 $\det(\rho)$，对应方式由 $n=1$ 的对应给出。
  \item 与扭相容，即若 $\chi$ 是 $k^\times$ 的特征，则
    $\rho_{\pi\otimes\chi}=\rho\otimes\chi$。
  \item $\pi_\rho$ 平方可积，当且仅当 $\rho(W'_k)$ 不落在 $\GL_n(\mathbb{C})$ 的
    真 Levi 子群中。
  \item $\pi_\rho$ 缓增，当且仅当 $\rho(W_k)$ 有界。
  \item 若 $H$ 为 $k$ 上连通约化群，$H(k)\to\GL_n(k)$ 是核与余核均交换的 $k$ 态射，
    则 $\GL_n(k)$ 的这些双射与 $H(k)$ 的类似映射之间须有相容性。
\end{enumerate}
\par\medskip
\begin{annote}
平方可积判据中的 $\rho(W'_k)$ 指整个 Weil--Deligne 表示的像，含 $N$ 的信息。
以 $\GL_2$ 的 Steinberg 表示为例，对应 $\rho'=\omega\otimes\mathrm{Sp}(2)$：
$\rho(W_k)$ 的像是形如
$\mathrm{diag}(\|w\|^{1/2},\|w\|^{-1/2})$ 的对角 torus，它本身就是一个真 Levi
子群；但 $W'_k$ 的像还含 $N$，其 Zariski 闭包落在 Borel 子群中而不落在任何真 Levi
子群中，故判据给出平方可积，与事实相符。若只看 $\rho(W_k)$，就会把 Steinberg 表示
误判为非平方可积。
\par\vspace{0.15em}{\raggedleft 【来源：Borel, \textit{Automorphic $L$-functions}, Proc.\ Symp.\ Pure Math.\ 33 (1979), Part 2, 27--61；Zelevinsky, Ann.\ Sci.\ \'Ec.\ Norm.\ Sup.\ 13 (1980), 165--210（A）】\par}
\end{annote}

关于相容性条件 (vi) 的更多细节，参看 Borel 在 Corvallis 的文章 [4] 或本书中配套的
文章 [14]。这与 Langlands 的一般函子性猜想有关。Langlands 本人从未把猜想的这一形式
与他的函子性一般原理分开 [40]。注意这一表述完全没有提到 $L$ 函数。

对非阿基米德局部域，先是在 Jacquet 与 Langlands 的书 [31] 中，接着在 Deligne 的
工作 [18] 中，局部 Langlands 猜想的一种表述不再用表示论性质，而用所论两集合的复解析
不变量即 $L$ 函数来陈述。Deligne 对 $\GL_2$ 给出了完整表述。正是在这篇文章中，他首次
使用了他在 [19] 中于 $\ell$ 进表示背景下引入的 Weil--Deligne 群，以便在非阿基米德
局部域上 $\GL_2$ 的情形得到正确的表述。

\par\medskip\noindent\textbf{局部 Langlands 猜想 II.}\ 设 $k$ 为非阿基米德局部域。
对每个 $n\ge1$，存在双射 $\mathcal{A}_n(k)\to\operatorname{Rep}_n(W'_k)$，
记作 $\pi\mapsto\rho_\pi$，具有以下性质。
\begin{enumerate}[label=(\roman*),leftmargin=*]
  \item 当 $n=1$ 时，该双射由局部类域论给出，并归一化使 $k$ 的一致化子对应几何
    Frobenius 元。
  \item 对任意 $\pi\in\mathcal{A}_n(k)$ 与 $\pi'\in\mathcal{A}_{n'}(k)$，有
    \[
      L(s,\rho_\pi\otimes\rho_{\pi'})=L(s,\pi\times\pi'),\qquad
      \varepsilon(s,\rho_\pi\otimes\rho_{\pi'},\psi)
      =\varepsilon(s,\pi\times\pi',\psi).
    \]
  \item 对任意 $\pi\in\mathcal{A}_n(k)$，$\rho_\pi$ 的行列式对应 $\pi$ 的中心特征，
    对应方式由局部类域论给出。
  \item 对任意 $\pi\in\mathcal{A}_n(k)$，有 $\rho_{\tilde\pi}=\tilde\rho_\pi$。
  \item 对任意 $\pi\in\mathcal{A}_n(k)$ 与 $k^\times$ 的任意有限阶特征 $\chi$，有
    $\rho_{\pi\otimes\chi}=\rho_\pi\otimes\chi$。
\end{enumerate}
\par\medskip
\begin{annote}
归一化是整条对应的定标点：局部类域论把一致化子 $\varpi$ 送到几何 Frobenius
$\Phi$，于是 $\det\rho_\pi$ 与中心特征 $\omega_\pi$ 在 $\varpi$ 上取值相同。若改用
算术 Frobenius（即 $\Phi^{-1}$），则每个 $\rho_\pi$ 都要乘上一个 $\omega^s$ 型的
特征，所有 $\varepsilon(s,\pi\times\pi',\psi)$ 与 $L$ 函数都随之整体扭转，定理 3.2、
3.3 的陈述也要改写。Harris--Taylor 与 Henniart 的论文都反复申明这一归一化，原因即在
于此。
\par\vspace{0.15em}{\raggedleft 【来源：Harris--Taylor, \textit{The geometry and cohomology of certain simple Shimura varieties}, Ann.\ of Math.\ Studies 151 (2001), Introduction；Henniart, Invent.\ math.\ 139 (2000), 439--455（A）】\par}
\end{annote}

对这些猜想所提供的图景，有两种思考方式。若把信息的主要传递方向看作从
$\mathcal{A}_n(k)$ 到 $\operatorname{Rep}_n(W'_k)$，则可视为 Langlands 对非 Abel
局部类域论的表述。若把信息的主要传递方向看作从 $\operatorname{Rep}_n(W'_k)$ 到
$\mathcal{A}_n(k)$，则给出 $\GL_n(k)$ 的不可约容许表示的算术参数化。这就是
$\mathcal{A}_n(k)$ 的算术 Langlands 分类。

\subsection{\texorpdfstring{$k$}{k} 为阿基米德局部域，即 \texorpdfstring{$k=\mathbb{R}$ 或 $\mathbb{C}$}{k=R 或 C}}

此时 $\mathcal{G}_k$ 已被充分理解：它或为 $\mathbb{Z}/2\mathbb{Z}$，或为平凡群。
所以此时信息的传递方向相反。Langlands 于 1973 年前后以极大的普遍性完成了这件事 [38]，
不仅对 $\GL_n$，而且对一般实约化群。阿基米德局部域没有 Weil--Deligne 群。表示论形式
就是如今所称的实群表示的 Langlands 分类或 Langlands 参数。事实上，Langlands 是把这件事
与 $\GL_n$ 情形用 $\operatorname{Rep}_n(W_k)$（一般 $G$ 情形用容许同态
$W_k\to{}^LG$）给出的算术参数化结合起来做的。深刻而有趣的部分，是用这些映射所提供的
信息对表示进行分类，特别是它们与离散序列构造的关系。

\par\medskip\noindent\textbf{定理 3.1.}\ 设 $k$ 为 $\mathbb{R}$ 或 $\mathbb{C}$。
则存在一族自然的双射
\[
  \operatorname{Rep}_n(W_k)\;\longleftrightarrow\;\mathcal{A}_n(k),
\]
满足局部 Langlands 猜想版本 I 的性质 (i)--(vi)。
\par\medskip

与实代数群通常的 Langlands 分类的确切关系，见 [38]。所证明的陈述当然就是 Langlands
最初给出的陈述，而这很可能正是猜想中那些精确条件的动因。再次注意，这些条件是表示论的，
$L$ 函数与 $\varepsilon$ 因子不起作用。

\subsection{\texorpdfstring{$k$}{k} 为非阿基米德局部域}

近年来，局部 Langlands 猜想的这第二种形式已在非阿基米德局部域上得到确立：先是
Laumon、Rapoport 与 Stuhler 于 1993 年处理正特征情形 [44]，随后 Harris 与 Taylor
于 1999 年 [27]、Henniart 于 2000 年 [30] 处理特征零情形。两种情形下，对应都是在
不可约 Galois 表示与超尖点表示之间的对应这一层面上建立的。原猜想中许多表示论上的预期
性质，已被扭 $L$ 函数的相等，即相应复解析不变量族的相等，所取代。

以 $\mathcal{A}_n^0(k)_f$ 记 $\GL_n(k)$ 的中心特征为有限阶的不可约容许表示的同构类
之集。于是 Laumon、Rapoport 与 Stuhler 的定理如下 [44]。

\par\medskip\noindent\textbf{定理 3.2.}\ 设 $k$ 为特征 $p>0$ 的局部域。对每个
$n\ge1$，存在双射
$\mathcal{A}_n^0(k)_f\to\operatorname{Rep}_n^0(\mathcal{G}_k)$，记作
$\pi\mapsto\rho_\pi$，满足局部 Langlands 猜想版本 II 的条件 (i)--(v)。
\par\medskip

当局部域 $k$ 的特征为 $0$ 时，Harris 与 Taylor [27] 以及 Henniart [28] 所确立的
局部 Langlands 猜想有完全相同的陈述。

\par\medskip\noindent\textbf{定理 3.3.}\ 设 $k$ 为特征 $0$ 的局部域。对每个
$n\ge1$，存在双射
$\mathcal{A}_n^0(k)_f\to\operatorname{Rep}_n^0(\mathcal{G}_k)$，记作
$\pi\mapsto\rho_\pi$，满足局部 Langlands 猜想版本 II 的条件 (i)--(v)。
\par\medskip

这些证明在 Galois 侧与自守侧都使用 $\Ql$ 表示，再通过 $\Ql$ 与 $\mathbb{C}$ 之间的
同构 $\iota$，翻译成上述用复解析 $L$ 函数表述的陈述。

\par\medskip\noindent\textbf{约化与构造。}\ 在任何非阿基米德局部情形，证明都经过一串
相同的约化，把问题归结为证明存在一个具有所需性质的映射。三个证明的区别在于用来证明
至少存在一个对应的构造。

约化本质上分三步。设已有任一对应
$\mathcal{A}_n^0(k)_f\to\operatorname{Rep}_n^0(\mathcal{G}_k)$，仍记作
$\pi\mapsto\rho_\pi$，且满足定理的 (i)--(v)。

1.~\textbf{单射性：$L$ 函数的极点。}\ 对
$\rho,\rho'\in\operatorname{Rep}_n^0(\mathcal{G}_k)$，$L(s,\rho\otimes\rho')$
在 $s=0$ 有极点当且仅当 $\rho'\cong\tilde\rho$。类似地，若 $\pi,\pi'$ 都在
$\mathcal{A}_n^0(k)$ 中，则 $L(s,\pi\times\pi')$ 在 $s=0$ 有极点当且仅当
$\pi'\cong\tilde\pi$。于是，任何满足 (ii) 的这种对应自动是单射。

2.~\textbf{双射性：数值局部 Langlands。}\ 对
$\rho\in\operatorname{Rep}_n^0(\mathcal{G}_k)$，以 $a(\rho)$ 记 $\rho$ 的
Artin 导子的指数 [51]，它由 $\varepsilon$ 因子 $\varepsilon(s,\rho,\psi)$ 决定。
把 $\mu\colon k^\times\to\mathbb{C}^\times$ 经局部类域论等同于 Galois 群的一个特征。
以 $\operatorname{Rep}_n^0(\mathcal{G}_k)_{m,\mu}$ 记
$\operatorname{Rep}_n^0(\mathcal{G}_k)$ 中满足 $a(\rho)=m$ 与
$\det(\rho)=\mu$ 的 $\rho$ 之集，则此集有限。类似地，以
$\mathcal{A}_n^0(k)_{m,\mu}$ 记 $\mathcal{A}_n^0(k)$ 中满足 $f(\pi)=m$ 与中心特征
$\omega_\pi=\mu$ 的 $\pi$ 之集，其中 $f(\pi)$ 是 Jacquet、Piatetski-Shapiro 与
Shalika 的导子的指数 [13]，此集也有限。数值局部 Langlands 猜想的陈述是：对固定的
$m\in\mathbb{Z}_+$ 与有限阶乘法特征 $\mu$，有
\[
  \bigl|\operatorname{Rep}_n^0(\mathcal{G}_k)_{m,\mu}\bigr|
  =\bigl|\mathcal{A}_n^0(k)_{m,\mu}\bigr|.
\]
该猜想已由 Henniart 于 1988 年确立 [30]。由于 (ii) 与 (iii) 保证我们的对应保持特征
$\mu$ 与导子，一旦知道对应是单射，数值局部 Langlands 猜想便给出它是满射，从而是双射。
\begin{annote}
此处两处引文编号与原书参考文献表不符，应作互换。(a)~数值局部 Langlands 猜想是
Henniart 1988 年的结果，即 [28] \textit{La conjecture de Langlands num\'erique
pour $\GL(n)$}, Ann.\ Sci.\ \'ENS 21 (1988), 497--544；正文标作 [30]，而 [30] 是
2000 年的简化证明。(b)~上段「Harris and Taylor [27] and Henniart [28]」中的 [28]
应为 [30]，因为定理 3.3 的特征零情形正是 2000 年那篇。译文照排原编号，读者按年份核对
即可。
\par\vspace{0.15em}{\raggedleft 【来源：原书参考文献表 [28]、[30]；Henniart, Invent.\ math.\ 139 (2000), 439--455（A）】\par}
\end{annote}

3.~\textbf{唯一性：局部逆定理。}\ 满足 (i)--(v) 的对应的唯一性，是 $\GL(n)$ 局部
逆定理的推论。该结果最早由 Jacquet、Piatetski-Shapiro 与 Shalika 陈述 [32]，
但第一个正式发表的证明由 Henniart 给出 [29]，其动机正是这一应用。陈述如下。
设 $\pi,\pi'$ 都属于 $\mathcal{A}_n^0(k)$，且扭 $\varepsilon$ 因子相同，即对一切
$1\le m\le n-1$ 与一切 $\tau\in\mathcal{A}_m^0(k)$ 有
\[
  \varepsilon(s,\pi\times\tau,\psi)=\varepsilon(s,\pi'\times\tau,\psi),
\]
则 $\pi\cong\pi'$。（注意相应的扭 $L$ 函数都恒等于 $1$ [13]。）由此对 $n$ 归纳可知，
任何这样的（此时已是双射的）满足 (i)--(v) 的对应必定唯一。
\begin{annote}
「相应的扭 $L$ 函数都恒等于 $1$」的理由是：$\pi$ 与 $\tau$ 都超尖点，且
$m\le n-1$，故 $m\neq n$。Jacquet--Shalika 的理论给出，两个超尖点表示的
Rankin--Selberg 卷积 $L(s,\pi\times\tau)$ 或恒等于 $1$，或（仅当 $m=n$ 且
$\tau\cong\tilde\pi$ 时）出现极点。既然这里 $m<n$，它只能恒等于 $1$。于是逆定理实际
只检验 $\varepsilon$ 因子，这也是它被称为「$\varepsilon$ 因子逆定理」的由来。
\par\vspace{0.15em}{\raggedleft 【来源：Jacquet--Piatetski-Shapiro--Shalika, C.~R.\ Acad.\ Sci.\ Paris 289 (1979), A59--A61；Cogdell, 本书第 9 章 \S4（A）】\par}
\end{annote}

这三步把局部 Langlands 猜想归结为存在某个满足 (i)--(v) 的对应的问题。正是这一存在性
问题，在正特征情形由 Laumon、Rapoport 与 Stuhler 解决，在特征零情形由 Harris 与
Taylor 继而 Henniart 解决。

4.~\textbf{存在性：整体几何构造。}\ 在所有情形，局部存在性都建立在确证 Galois 表示
与自守表示之间整体对应的某些实例之上。注意此时 Galois 侧与自守侧用的都是 $\ell$ 进表示。

对特征 $p>0$ 的 $k$，Laumon、Rapoport 与 Stuhler 从局部表示
$\pi\in\mathcal{A}_n^0(k)$ 出发。他们把 $k$ 实现为特征 $p$ 的整体域 $K$ 的一个
局部分量，即 $k=K_v$（$v$ 为 $K$ 的某位点），再把 $\pi$ 嵌入为秩 $n$ 的整体除代数
$D$ 的阿代尔点 $D^\times(\mathbb{A})$ 上尖点表示 $\Pi$ 在 $v$ 处的局部分量，使得
$D^\times(K_v)=\GL_n(k)$ 且 $\Pi_v=\pi$。他们在 $\mathcal{D}$-椭圆模（$\mathcal{D}$
为 $D$ 中的一个序）的模空间的 $\ell$ 进上同调上整体地实现
$\mathcal{G}_K\times D^\times(\mathbb{A})$ 的作用，使得该上同调的分解中出现
$\mathcal{G}_K\times D^\times(\mathbb{A})$ 的表示 $R\otimes\Pi$。按构造
$\Pi_v=\pi$，他们取 $R_v=\rho_\pi$。由其构造的性质，他们能验证 (i)--(v) 成立。
于是他们通过一个整体几何构造与一个受限的整体对应，建立了所需的局部存在性陈述。

Harris 与 Taylor 对特征零非阿基米德域的局部 Langlands 猜想的证明，在精神上与
Laumon、Rapoport 与 Stuhler 的特征 $p$ 证明相似。他们把 $\mathcal{D}$-椭圆模的
模空间换成与 Kottwitz 的酉群 $U_n$ 相联系的某些「简单 Shimura 簇」。他们把 $k$ 实现
为数域 $K$ 的一个局部分量，即 $k=K_v$（$v$ 为 $K$ 的某位点），再把 $\pi$ 嵌入为
某个（扭）酉群的秩 $n$ 尖点表示 $\Pi$ 在 $v$ 处的局部分量。然后他们在相应 Shimura 簇
的 $\ell$ 进上同调中，实现这些尖点表示与整体 Galois 表示之间的整体对应。通过在一个
坏约化位点处局部地研究所得的对应，他们找到这样一个局部表示：$\GL_n(k)\times
D_n^\times(k)\times W_k$ 在其上有不可约的作用，作用由
$\pi\otimes JL(\pi)\otimes\rho_\pi$ 给出，其中 $D$ 是 $k$ 上秩 $n$、Hasse 不变量
为 $1/n$ 的除代数，$JL(\pi)$ 是 $\pi$ 在局部 Jacquet--Langlands 对应下的像，而
$\rho_\pi$ 由此得到定义。同样，由其构造他们能验证该对应满足条件 (i)--(v)。注意他们
不仅得到局部 Langlands 对应的几何实现，还同时得到局部 Jacquet--Langlands 对应的
实现。
\begin{annote}
这里同时出现了两个对应。$JL$ 是局部 Jacquet--Langlands 对应：把 $\GL_n(k)$ 的
超尖点表示送到 $k$ 上秩 $n$、Hasse 不变量 $1/n$ 的除代数 $D$ 的乘法群
$D^\times(k)$ 的表示，并保持中心特征与 $\varepsilon$ 因子。构造的妙处在于，简单
Shimura 簇在坏约化处的上同调同时带有 $\GL_n(k)$ 与 $D^\times(k)$ 的作用，于是
$\pi\otimes JL(\pi)$ 自动出现，$\rho_\pi$ 作为第三个因子的作用被「读出来」。这也
解释了为何该证明一举给出局部 Langlands 与局部 Jacquet--Langlands 两个对应。
\par\vspace{0.15em}{\raggedleft 【来源：Harris--Taylor 2001, Introduction 与 Ch.~1；Carayol, S\'eminaire Bourbaki 857, Ast\'erisque 266 (2000), 191--243（A）】\par}
\end{annote}

Henniart 对特征零非阿基米德域的局部 Langlands 猜想的证明，同样使用整体构造，
但用法远不那么沉重。特别地，他不给出对应的几何实现。对 Henniart 而言，陈述与证明都
最自然地给出从 Galois 侧到自守侧的双射 $\rho\mapsto\pi_\rho$。Henniart 从 $\mathcal{G}_k$
的一个行列式为有限阶的不可约表示 $\rho$ 出发。它可分解为某个有限扩张 $F/k$ 的
$\operatorname{Gal}(F/k)$ 的表示。利用 Brauer 归纳，他把 $\rho$ 写成单项表示之和。
这些特征可经局部类域论提升到自守侧，于是他须证明相应的自守诱导表示之和存在且超尖点。
所得的尖点表示就是 $\pi_\rho$。这一步他仍是把局部情形嵌入整体情形，再引用 Harris
早些时候建立的某些整体自守诱导的弱情形 [25]。Harris 的结果依赖基变换理论，以及
Clozel 关于把 $\ell$ 进表示附于 $\GL_n(\mathbb{A})$ 的自守表示的工作 [12]，而后者
又依赖 Kottwitz 关于某些酉 Shimura 簇良好约化的工作。所以归根到底仍有一个整体几何
构造，只是比 Harris 与 Taylor 所用的简单。Henniart 的证明更多地使用 $L$ 函数，更少
地使用几何。他的证明更短、更解析，但不给出对应的几何实现。

关于这些结果更完整的提要，可见 Carayol 的 Séminaire Bourbaki 报告 [9, 10]。

\subsection{补充}

为完成局部 Langlands 对应，须在 Galois 侧考虑 Weil--Deligne 群的所有适当表示，
在自守侧考虑 $\GL_n(k)$ 的所有不可约容许表示。为此，第一步是去掉中心特征有限阶的
条件。这只需在对应的 Galois 侧把 Galois 群换成 Weil 群；自守侧仍是超尖点表示。
然后，要过渡到 $\GL_n(k)$ 的所有容许表示，就要用 Weil--Deligne 群的表示。在表示论上，
Galois 侧从 Weil 群表示到 Weil--Deligne 群表示的过渡，恰好镜像自守侧从超尖点表示到
不可约表示的过渡，这一点由 Bernstein 与 Zelevinsky 证明（见 [1, 60]，特别是 [60]
的第 10 节，或 [47]）。于是由 Laumon、Rapoport、Stuhler、Harris、Taylor 与 Henniart
的结果，完整的局部 Langlands 对应随之成立。

尽管有这些结果，局部 Langlands 猜想的研究仍在继续。上述证明给出对应的存在性，某些情形
还提供显式的几何模型。它给出若干不变量的匹配，例如导子与局部 $L$ 函数。然而，为了应用，
人们希望有局部 Langlands 对应的显式版本，特别是用 Bushnell--Kutzko 紧诱导数据 [6]
表述的 $\GL_n$ 超尖点表示的显式版本。显式局部 Langlands 对应的探索目前由 Bushnell、
Henniart、Kutzko 等人推进。

\section{整体 Langlands 猜想}

与局部情形一样，整体 Langlands 猜想最基本的形式，应是（Abel）整体类域论的非 Abel
推广。当 Deligne 指出在局部非阿基米德情形必须引入 Weil--Deligne 群时，人们随即意识到
似乎并不存在自然的整体 Weil--Deligne 群。这促使人们寻找一个整体群来替代
Weil--Deligne 群。这正是 Langlands 文章 [39] 的目的之一。如今人们相信，这个被
Ramakrishnan 称为猜想中的 Langlands 群 $\mathcal{L}_k$ 的群，应与同样属于猜想的
$k$ 的动机 Galois 群 $\mathcal{M}_k$ 有关 [46]。

\subsection{\texorpdfstring{$k$ 为特征 $p>0$ 的整体域}{k 为特征 p>0 的整体域}}

尽管有这些困难，Drinfeld 仍对整体函数域表述并证明了整体 Langlands 猜想的一个版本 [23]，
它把 Galois 群本身的不可约二维表示与 $\GL_2(\mathbb{A})$ 的不可约尖点表示联系起来。
这是 Laumon、Rapoport 与 Stuhler 局部定理的整体类比，而那里并不需要 Weil--Deligne 群。
应当强调，Drinfeld 的结果是 1970 年代得到的，比 Laumon、Rapoport 与 Stuhler 的工作
早了好几年。近年来，Drinfeld 的工作被 L.~Lafforgue 推广，给出了函数域上 $\GL_n$ 的整体
Langlands 猜想的证明 [35]。

Drinfeld 与 Lafforgue 所建立的整体 Langlands 猜想的表述，与上述非阿基米德局部情形
本质相同，只有几处改动，我们下面加以说明。取 $k$ 为特征 $p$ 的有限域 $\mathbb{F}$ 上
光滑、射影、几何连通曲线 $X$ 的函数域。固定与 $p$ 不同的素数 $\ell$，以及同构
$\iota\colon\Ql\to\mathbb{C}$。

在 Galois 侧，他们考虑如下不可约连续 $\ell$ 进表示 $\rho\colon\mathcal{G}_k\to
\GL_n(\Ql)$ 的同构类：在有限多个位点之外无分歧（如第 1 节所述），且行列式为有限阶。
我们把这些记作 $\operatorname{Rep}_n^0(\mathcal{G}_k;\Ql)_f$。在自守侧，他们考虑
中心特征为有限阶的 $\Ql$ 值尖点表示空间 $\mathcal{A}_n^0(k;\Ql)_f$，如第 2 节所述。
类比局部情形，整体 Langlands 猜想的一个合理表述如下。

\par\medskip\noindent\textbf{特征 $p$ 的整体 Langlands 猜想.}\ 对每个 $n\ge1$，
存在双射 $\mathcal{A}_n^0(k;\Ql)_f\to\operatorname{Rep}_n^0(\mathcal{G}_k;\Ql)_f$，
记作 $\pi\mapsto\rho_\pi$，具有以下性质。
\begin{enumerate}[label=(\roman*),leftmargin=*]
  \item 当 $n=1$ 时，该双射由整体类域论给出。
  \item 对任意 $\pi\in\mathcal{A}_n^0(k;\Ql)_f$ 与
    $\pi'\in\mathcal{A}_{n'}^0(k;\Ql)_f$，有
    \[
      L(s,\rho_\pi\otimes\rho_{\pi'})=L(s,\pi\times\pi'),\qquad
      \varepsilon(s,\rho_\pi\otimes\rho_{\pi'})=\varepsilon(s,\pi\times\pi').
    \]
  \item 对任意 $\pi\in\mathcal{A}_n^0(k;\Ql)_f$，$\rho_\pi$ 的行列式对应 $\pi$ 的
    中心特征，对应方式由整体类域论给出。
  \item 对任意 $\pi\in\mathcal{A}_n^0(k;\Ql)_f$，有
    $\rho_{\tilde\pi}=\tilde\rho_\pi$。
  \item 对任意 $\pi\in\mathcal{A}_n^0(k;\Ql)_f$ 与 $k^\times$ 的任意有限阶特征
    $\chi$，有 $\rho_{\pi\otimes\chi}=\rho_\pi\otimes\chi$。
  \item 这些整体双射应与局部 Langlands 猜想的局部双射相容。
\end{enumerate}
\par\medskip

若取尖点表示 $\pi\in\mathcal{A}_n^0(k;\Ql)_f$，以 $S(\pi)$ 记一个有限位点集
$S(\pi)\subset|X|$，使得对一切 $x\notin S(\pi)$，$\pi_x$ 都无分歧。对
$x\notin S(\pi)$，局部表示 $\pi_x$ 完全由 $\GL_n(\Ql)$ 中的一个半单共轭类
$A_{\pi_x}$ 决定；经 $\iota$ 把它与 $\GL_n(\mathbb{C})$ 等同，称之为 $\pi_x$ 的
Satake 类或 Satake 参数 [4]。由于该参数决定且被决定为 $x$ 处相应无分歧 Hecke 代数
$\mathcal{H}_x$ 的一个特征 [50]，$A_{\pi_x}$ 的特征值
$z_1(\pi_x),\dots,z_n(\pi_x)$ 也称为 $\pi_x$ 的 Hecke 特征值。这些 Hecke 特征值
完全决定 $\pi_x$。再由 $\GL_n$ 的强重数一定理 [13]，几乎所有
$x\notin S(\pi)$ 的 Hecke 特征值组
$\{z_1(\pi_x),\dots,z_n(\pi_x)\}$ 完全决定 $\pi$。

若取 Galois 表示 $\rho\in\operatorname{Rep}_n^0(\mathcal{G}_k;\Ql)_f$，则也有
有限位点集 $S(\rho)$，使 $\rho$ 在一切 $x\notin S(\rho)$ 处无分歧。对
$x\notin S(\rho)$，$x$ 处几何 Frobenius 元 $\Phi_x$ 的像
$\rho(\Phi_x)$ 是 $\GL_n(\Ql)\cong\GL_n(\mathbb{C})$ 中一个确定的半单共轭类。
$\rho(\Phi_x)$ 的特征值 $z_1(\rho_x),\dots,z_n(\rho_x)$ 称为 $\rho_x$ 的
Frobenius 特征值。这些 Frobenius 特征值完全决定 $\rho_x$；由 Chebotarev 密度定理，
几乎所有 $x\notin S(\rho)$ 的 Frobenius 特征值组
$\{z_1(\rho_x),\dots,z_n(\rho_x)\}$ 完全决定 $\rho$ 本身。

Drinfeld 对 $n=2$ [23]、Lafforgue 对 $n\ge3$ [35] 所建立的结果如下；我们将在下面
说明，它与上述陈述等价。

\par\medskip\noindent\textbf{定理 4.1.}\ 设 $k$ 为如上特征 $p$ 的整体函数域。
对每个 $n\ge1$，存在唯一的双射
$\mathcal{A}_n^0(k;\Ql)_f\to\operatorname{Rep}_n^0(\mathcal{G}_k;\Ql)_f$，
记作 $\pi\mapsto\rho_\pi$，使得对每个尖点表示
$\pi\in\mathcal{A}_n^0(k;\Ql)_f$，都有 Hecke 特征值与 Frobenius 特征值的相等
\[
  \{z_1(\pi_x),\dots,z_n(\pi_x)\}
  =\{z_1(\rho_{\pi,x}),\dots,z_n(\rho_{\pi,x})\},
\]
对一切 $x\notin S$ 成立，其中 $S$ 是包含 $S(\pi)\cup S(\rho_\pi)$ 的有限位点集；
等价地，有部分复解析 $L$ 函数的相等
\[
  L^S(s,\pi)=L^S(s,\rho_\pi).
\]
\par\medskip

为看出这确实给出上述猜想的陈述，首先注意：对一切 $x\notin S$，相应 Hecke 特征值与
Frobenius 特征值在这些位点处的相等，与未分歧位点处的局部 Langlands 猜想一致。
故 (vi) 在这些位点得到满足。其次，对
$\pi\in\mathcal{A}_n^0(k;\Ql)_f$，$L$ 函数有 Euler 乘积分解
$L(s,\pi)=\prod_{x\in|X|}L(s,\pi_x)$，且对一切 $x\notin S(\pi)$，局部 $L$ 函数为
$L(s,\pi_x)=\det(1-\iota(A_{\pi_x})q^{-s})^{-1}$ [13]。于是
$\pi_x$（$x\notin S(\pi)$）的 Hecke 特征值由这些位点处的局部 $L$ 因子决定，反之
亦然。类似地，对 $\rho\in\operatorname{Rep}_n^0(\mathcal{G}_k;\Ql)_f$，整体
$L$ 函数同样有分解 $L(s,\rho)=\prod_{x\in|X|}L(s,\rho_x)$，只是现在
$L(s,\rho_x)=\det(1-\iota(\rho(\Phi_x))q_x^{-s})^{-1}$。于是
$\rho(\Phi_x)$（$x\notin S(\rho)$）的局部 Frobenius 特征值由这些位点处的局部
$L$ 因子决定，反之亦然。因此对包含 $S(\pi)\cup S(\rho_\pi)$ 的有限集 $S$，
确实有部分 $L$ 函数的相等 $L^S(s,\pi)=L^S(s,\rho_\pi)$，如所述。对 $L(s,\pi)$
与 $L(s,\rho_\pi)$ 同时使用整体函数方程，并利用整体 $\varepsilon$ 因子的局部分解，
标准的 $L$ 函数技巧给出：事实上 $S(\pi)=S(\rho_\pi)$，在这些位点处
$L(s,\pi_x)=L(s,\rho_{\pi,x})$，并且一般而言 $\rho_\pi$ 在 $x$ 处对 $k$ 的局部
Galois 群的限制，在局部 Langlands 猜想下对应 $\pi_x$。故 (vi) 一般地成立；由此，
(ii)--(v) 可由 Euler 乘积分解以及局部猜想的相应陈述推出。

\par\medskip\noindent\textbf{约化与构造。}\ 与局部情形一样，证明经过一串约化，把问题
归结为证明存在一个具有所需性质的映射。存在性随后由整体构造建立，用的是某个
$\GL_n$ 在其上作用的模概形的上同调。

1.~\textbf{唯一性与双射性。}\ 给定如上所述的一个整体双射，其唯一性早已知道可由自守侧
的强重数一定理 [13] 与 Galois 侧的 Chebotarev 密度定理推出。这些强唯一性原理还蕴涵：
任何满足定理条件的映射 $\pi\mapsto\rho_\pi$ 与 $\rho\mapsto\pi_\rho$ 必定互为逆双射。

2.~\textbf{Piatetski-Shapiro 与 Deligne 的归纳程序。}\ 这一归纳原理由 Deligne 于
1980 年在 IHES 的讨论班上讲出，后来记录于 [41]。它把定理的证明归结为下述看似更弱的
存在性陈述。

\par\medskip\noindent\textbf{定理 4.2.}\ 对每个 $n\ge1$，存在映射
$\mathcal{A}_n^0(k;\Ql)_f\to\operatorname{Rep}_n^0(\mathcal{G}_k;\Ql)_f$，
记作 $\pi\mapsto\rho=\rho_\pi$，使得 Hecke 特征值与 Frobenius 特征值相等
\[
  \{z_1(\pi_x),\dots,z_n(\pi_x)\}
  =\{z_1(\rho_{\pi,x}),\dots,z_n(\rho_{\pi,x})\},
\]
对几乎所有 $x\notin S(\pi)\cup S(\rho_\pi)$ 成立。
\par\medskip

事实上，设已对 $r=1,\dots,n-1$ 建立了映射 $\pi\mapsto\rho_\pi$
（$\pi\in\mathcal{A}_r^0(k;\Ql)_f$）。则利用 Grothendieck 的整体函数方程、
整体 Galois $\varepsilon$ 因子的分解 [41]，以及 $\GL_n$ 的逆定理 [13, 15, 16]，
便可自动得到逆映射 $\rho\mapsto\pi_\rho$
（$\rho\in\operatorname{Rep}_r^0(\mathcal{G}_k;\Ql)_f$，$r=1,\dots,n$）。

3.~\textbf{存在性（[23, 35]）。}\ 自 Weil 以来人们就知道，函数域上 $\GL_n$ 有自然的
模问题：曲线 $X$ 上秩 $n$ 向量丛的同构类由双陪集
$\GL_n(k)\backslash\GL_n(\mathbb{A})/\GL_n(\mathcal{O})$ 参数化。要得到所需的映射，
还需要一点额外的结构，于是 Drinfeld 继而 Lafforgue 考虑了秩 $n$ shtuka（带级结构）
的（紧化）Deligne--Mumford 叠 $\mathcal{V}$，它其实是 $X\times X$ 上的叠。整体
Hecke 代数 $\mathcal{H}$ 通过对应自然地作用在此叠上，相应的 $\ell$ 进上同调
$H^*_c(\bar k\otimes_k\mathcal{V};\Ql)$ 于是承载 $\mathcal{H}$ 与
$\mathcal{G}_k\times\mathcal{G}_k$ 的同时表示。接着用几何 Grothendieck--Lefschetz
迹公式计算此表示的迹，再与 Arthur--Selberg 迹公式的输出比较，证明
$\mathcal{H}\times\mathcal{G}_k\times\mathcal{G}_k$ 的导出表示
$\pi\otimes\rho_\pi\otimes\tilde\rho_\pi$ 确实出现在该上同调中。这一构造是归纳的，
并且实质上用到了所有东西。
\begin{annote}
双陪集 $\GL_n(k)\backslash\GL_n(\mathbb{A})/\GL_n(\mathcal{O})$ 分类曲线 $X$ 上
秩 $n$ 向量丛的同构类，这是 Weil 的观察。但它只给出「点」，不足以让 Hecke 算子以
对应（correspondence）的方式作用。shtuka 给向量丛配上 Frobenius 结构，其两个「腿」
给出 $X\times X$ 上的叠，Hecke 对应才得以几何化，$\mathcal{H}$ 与
$\mathcal{G}_k\times\mathcal{G}_k$ 才能同时作用在上同调上。
\par\vspace{0.15em}{\raggedleft 【来源：Drinfeld, \textit{Cohomology of compactified modules of $F$-sheaves of rank 2}, J.\ Soviet Math.\ 46 (1989), 1789--1821；Lafforgue, Invent.\ math.\ 147 (2002), 1--241（A）】\par}
\end{annote}

\par\medskip\noindent\textbf{补充。}\ 此结果的一个直接推论是 $\GL_n$ 的
Ramanujan--Petersson 猜想。该猜想此前已由 Drinfeld 对 $n=2$ [22] 证明，并由
Lafforgue 对 $n\ge3$ 部分证明 [34]。完整的解答由整体 Langlands 猜想推出。

\par\medskip\noindent\textbf{定理 4.3.}\ 对每个 $\pi\in\mathcal{A}_n^0(k)_f$ 与
每个位点 $x\notin S(\pi)$，有 $|z_i(\pi_x)|=1$。
\par\medskip

此外，Lafforgue [35] 还能推出 Deligne 的下述猜想 [20]。

\par\medskip\noindent\textbf{定理 4.4.}\ 曲线上的每个不可约局部系统 $\rho$，若其
行列式为有限阶，则为权 $0$ 的纯系统；并且，Frobenius 特征值的对称多项式生成
$\mathbb{Q}$ 的有限扩张。
\par\medskip
\begin{annote}
「权 $0$ 的纯」是 Weil 猜想的语言：对未分歧位点 $x$，$\rho(\Phi_x)$ 的每个特征值
$\alpha$ 在任一复嵌入下满足 $|\iota(\alpha)|=1$。定理 4.3 是它的自守对应物：
$|z_i(\pi_x)|=1$ 即 Ramanujan--Petersson 猜想。定理 4.4 后半句「Frobenius 特征值的
对称多项式生成 $\mathbb{Q}$ 的有限扩张」说明这些特征值是代数数，从而相应 $L$ 函数的
取值落在 $\overline{\mathbb{Q}}$ 中；这正是能把 $\ell$ 进与复两种 $L$ 函数加以比较
的前提。
\par\vspace{0.15em}{\raggedleft 【来源：Deligne, \textit{La conjecture de Weil II}, Publ.\ Math.\ IHES 52 (1980), 137--252；Lafforgue, Invent.\ math.\ 147 (2002)（A）】\par}
\end{annote}

此外，Lafforgue 还能断定：在曲线上，不可约局部 $\ell$ 进系统的概念不依赖 $\ell$ 的
选取；他并验证了「几何 Langlands 对应」中的下降断言。关于这些结果更完整的提要，
可见 Laumon 的报告 [42, 43]。

\subsection{\texorpdfstring{$k$ 为特征 $0$ 的整体域}{k 为特征 0 的整体域}}

数域情形的一般性结果少得可怜。不过，这种整体对应确有一些相当引人注目的例子。

\par\medskip\noindent\textbf{一般猜想。}\ 回顾 $n=1$ 的情形：整体类域论给出
$\mathcal{G}_k$ 的连续特征与 $k^\times\backslash\mathbb{A}^\times$ 的有限阶特征之间
的典范双射。要得到 $k^\times\backslash\mathbb{A}^\times$ 的全部特征，须再次把
Galois 群换成整体 Weil 群 $W_k$。

对 $n\ge2$，类比局部 Langlands 猜想，我们需要 Weil--Deligne 群的整体类比物。可惜
并无这样的类比物可用。于是猜想改以猜想中的 Langlands 群 $\mathcal{L}_k$ [46] 来
表述。人们至多希望 $\mathcal{L}_k$ 进入一个正合列
\[
  1\longrightarrow\mathcal{L}_k^0\longrightarrow\mathcal{L}_k
  \longrightarrow\mathcal{G}_k\longrightarrow 1,
\]
其中 $\mathcal{L}_k^0$ 为复投射约化群。这应当进入一个交换图
\[
\begin{tikzcd}
  1 \arrow[r] & \mathcal{L}_k^0 \arrow[r] \arrow[d]
    & \mathcal{L}_k \arrow[r] \arrow[d]
    & \mathcal{G}_k \arrow[r] \arrow[d] & 1 \arrow[d] \\
  1 \arrow[r] & \mathcal{M}_k^0 \arrow[r]
    & \mathcal{M}_k \arrow[r]
    & \mathcal{G}_k \arrow[r] & 1
\end{tikzcd}
\]
其中 $\mathcal{M}_k$ 是同样属于猜想的动机 Galois 群 [46]。

按这些术语，一般地人们预期或猜想以下类型的整体对应 [11, 46]：
\begin{enumerate}[label=(\roman*),leftmargin=*]
  \item $\mathcal{G}_k$ 的不可约 $n$ 维表示，应与 $\GL_n(\mathbb{A})$ 的 Galois 型
    尖点表示一一对应。（这是对 $\pi_\infty$ 的限制。）
  \item $\mathcal{M}_k$ 的不可约 $n$ 维表示，应与 $\GL_n(\mathbb{A})$ 的代数尖点
    表示一一对应。它们是代数 Hecke 特征的类比物。
  \item $\mathcal{L}_k$ 的不可约 $n$ 维表示，应与 $\GL_n(\mathbb{A})$ 的全部尖点
    表示一一对应。
\end{enumerate}
当然，所有这些对应都应满足与局部猜想类似的性质，特别是保持 $L$ 因子与 $\varepsilon$
因子（带扭），以及与局部对应相容，等等。

实际上，真正一般性的结果所知极少。当前方法的一个困难似乎在于：$\GL_n$ 在数域上没有
自然的模问题。

\par\medskip\noindent\textbf{已知结果。}\ 若从自守侧出发去构造相应的 Galois 表示，
则有许多一般性的部分结果。

当 $n=2$、$k=\mathbb{Q}$ 时，有 Deligne 的基本结果 [17]，它基于 Eichler 与 Shimura
的奠基性工作，把每个对应于权 $\ge2$ 全纯新形式的 $\GL_2(\mathbb{A}_{\mathbb{Q}})$
尖点表示 $\pi$，附上一个相容的 $\ell$ 进表示系 $\rho=\rho_\pi$，使得
$L(s,\pi)=L(s,\rho)$。由此可推出这类形式的 Ramanujan--Petersson 猜想。Deligne 与
Serre [21] 在经典背景下把它推广到 $\mathbb{Q}$ 上的权 $1$ 形式。这些结果又被许多人
推广到全实域 $k$（仍是 $n=2$），其中包括 Rogawski--Tunnell [48]、Ohta [45]、
Carayol [8]、Wiles [58]、Taylor [52, 54] 与 Blasius--Rogawski [3]。对虚二次域，
有 Harris--Soudry--Taylor [26] 与 Taylor [53] 的工作。更完整的综述可参看 Blasius [2]
与 Taylor [55]。

对一般的 $\GL_n$ 与全实域 $k$，Clozel 能把相容的 $\ell$ 进表示系附于
$\GL_n(\mathbb{A}_k)$ 的尖点、代数、正则、自对偶表示，只要该表示在一两个有限位点处
具有特定类型的局部分量 [12]。

更为引人注目的是相反方向的结果，即从具体的 Galois 表示出发，证明它是模的。我们所指的
是 Langlands [37] 与 Tunnell [57]（以及 Taylor 等人的部分结果 [7, 56]）关于二次
复 Galois 表示模性的工作（强 Artin 猜想），以及 Wiles [59] 继而 Breuil、Conrad、
Diamond 与 Taylor [5] 关于 $\mathbb{Q}$ 上椭圆曲线（其 $\ell$ 进 Tate 模上的二维
Galois 表示）模性的工作。

\par\medskip
\noindent\textit{（原书本章末列参考文献 60 条，见原书 pp.~245--249，编号 [1]--[60]
与正文引用一一对应；条目为书目信息，本译本不逐条迻录，读者可按编号核对原书。）}
